% 付録G 演習の解答指針（巻末）
\section*{付録G\quad 演習の解答指針}
\addcontentsline{toc}{section}{付録G\quad 演習の解答指針}

各問 1--3 行の指針。詳細計算は読者に委ねる。

\paragraph{第1章}
(1) 放射優勢で $\mathcal H\propto a^{-1}$、物質優勢で $\mathcal H\propto a^{-1/2}$。
$\eta=\int c\,dx/\mathcal H\propto\int a^{1/2}da/a$ より物質優勢で $\eta\propto a^{1/2}$。
(2) $z_{\rm eq}=\Omega_{m0}/\Omega_{r0}-1$、$\Omega_{r0}$ は固定なので $z_{\rm eq}\propto\Omega_mh^2$。
(3) $\Omega_\Lambda\to0$ で後期の減速が増し $\eta_0$ は増える（NB1 で確認）。

\paragraph{第2章}
(1) $\Omega_{\gamma0}=a_{\rm rad}T_{\rm CMB}^4/(c^2\rho_{\rm crit,0})$ に数値代入で
$\Omega_{\gamma0}h^2\approx2.47\times10^{-5}$。
(2) $e^+e^-$ 消滅前後で $g_{*s}T^3a^3$ 保存、光子のみ加熱 $\Rightarrow(T_\nu/T_\gamma)^3=4/11$。
(3) 再結合期 $h\nu\sim k_BT_*\sim0.3$~eV $\ll m_ec^2=511$~keV なので Thomson 極限が有効。

\paragraph{第3章}
(1) Saha の $e^{-\epsilon_0/k_BT_b}$ と $n_b\propto\Omega_bh^2$ から、$\Omega_bh^2$ 増で
$z_*$ はわずかに増える（対数的）。
(2) $\tilde g=-\tau'e^{-\tau}$、変数変換 $d\tau=\tau'dx$ で
$\int\tilde g\,dx=\int_0^\infty e^{-\tau}d\tau=1$。
(3) $\Omega_b$ 増で再結合が早まり freeze-out 電離度はやや下がる（NB2）。

\paragraph{第4章}
(1) $\delta\rho\to\delta\rho-\dot{\bar\rho}\xi^0$、$\xi^0=\delta\rho/\dot{\bar\rho}$ を選べば
$\delta\rho=0$（一様密度スライス）。
(2) 計量\eqref{eq:metric} は非対角・非等方成分を持たず、スカラー2自由度を使い切るので
残留自由度はない。
(3) 付録A の対応表で $\Phi_{\rm book}=-\phi_{\rm MB}$ を確認。

\paragraph{第5章}
(1) 異方性応力はトレースレス $ij$ 成分から来て $\nabla^2$ 由来の $k^2$ で割られるため
$\Psi+\Phi\propto k^{-2}\sigma$。
(2) tight coupling で $\Theta_2\propto1/\tau'\to0$、式\eqref{eq:G2} で $\Psi\to-\Phi$。
(3) 式\eqref{eq:G1} の時間微分を落とし $ck/\mathcal H\gg1$ で $-c^2k^2\Phi/a^2=4\pi G\bar\rho\delta$。

\paragraph{第6章}
(1) 電子静止系の等方散乱を実験室系へ $v_b$ で1次ブーストすると $\hat p\cdot\mathbf v_b$、
多重極で $\tfrac13v_b$ が $\Theta_1$ 方程式に入る。
(2) 光子が得る運動量＝バリオンが失う運動量。$\rho_\gamma/\rho_b=4/(3R)$ 倍で符号反転 $\Rightarrow\tau'R(3\Theta_1+v_b)$。
(3) NB3 で $\Theta_\ell$ が高次へエネルギーを運ぶ（自由伝播）様子を観察。

\paragraph{第7章}
(1) 断熱・超地平線で $\Theta_0=-\Psi/2$（放射）。式の $\Theta_0=-\tfrac12\Psi$ に対応。
(2) $f_\nu\to0$ で $\Phi=-(1+\tfrac25 f_\nu)\Psi\to-\Psi$。
(3) $\Psi_{\rm ini}=-\tfrac23\mathcal R/(1+\tfrac4{15}f_\nu)$ に $\mathcal R=1$ を入れた値が
\S3.2 の $\Psi_{\rm ini}$。

\paragraph{第8章}
(1) 0次強結合 P2 の括弧 $[\Theta_1+v_b/3]=0\Rightarrow\Theta_1=-v_b/3$。
(2) 振動子\eqref{eq:osc} の強制項定数部が平衡点を $-(1+R)\Psi$ にずらす。
(3) NB7 で WKB の $\cos(kr_s)$ とフル数値の位相一致を確認。

\paragraph{第9章}
(1) $[1/k_D^2]=$ 長さ$^2$。$c^2/(\mathcal H^2|\tau'|)$ は長さ$^2$（$|\tau'|$ 無次元）。
(2) NB7 で $e^{-(k/k_D)^2}$ と数値振幅を重ね、$k\gtrsim0.2$ でずれ始める。
(3) $\ell_D=k_D\chi_*\approx0.149\times13890\approx2000$ 程度（減衰尾の落ち）。

\paragraph{第10章}
(1) $e^{ik\mu\Delta\eta}=\sum(2\ell+1)i^\ell P_\ell(\mu)j_\ell(k\Delta\eta)$ の $\ell=2$ 項。
(2) NB4 で SW/ISW/Doppler/四重極を個別プロット。
(3) 物質優勢・超地平線で $\Theta_0=-\tfrac12\Psi$ かつ $\Theta_0+\Psi=\tfrac13\Psi$
（重力赤方偏移 $+\Psi$ との和）。

\paragraph{第11章}
(1) $a_{\ell m}=4\pi(-i)^\ell\int\ldots$ の $4\pi$ は平面波の球面調和展開
$e^{i\mathbf k\cdot\hat n r}=4\pi\sum i^\ell j_\ell Y Y^*$ 由来。
(2) NB5 で $\mathcal D_\ell$ と Planck binned を重ね、$\ell_1\approx220$ を読む。
(3) $\sqrt{2/(2\cdot2+1)}=\sqrt{2/5}\approx63\%$。

\paragraph{第12章}
(1) $\ell_1=(1-0.27)\ell_A=(0.73)\times\vallA\approx\vallone$。
(2) NB5 でピーク検出、$H_2/H_1$ は $\Omega_bh^2$ 増で低下（圧縮増強）。
(3) late ISW は $\Lambda$ 優勢開始 $z\sim0.3$ の地平線、$\ell\lesssim30$。

\paragraph{第13章}
(1) NB6 で $z_{\rm re}$ を振ると $\tau\propto\int n_e d\eta$ が増え、$\ell=500$ 抑制
$=e^{-2\tau}-1$ が深くなる。
(2) $r_s\propto z_*^{-1/2}$ 程度、$\chi_*$ はほぼ不変 $\Rightarrow\partial\ln\ell_1/\partial\ln z_*>0$。
(3) $100\theta_*=$const 線上で $(\Omega_\Lambda,\Omega_k)$ を動かすと $\ell>30$ はほぼ不変。

\paragraph{第14章}
(1) tight coupling を切ると再結合前で stiff になり計算が遅く・不安定になる。
(2) NB8 で $n_k\in\{150,250,400\}$、中 $\ell$ の収束を確認。
(3) 細 $k$ 間隔を $2\times$ 粗くすると $|\Theta_\ell|^2$ の振動を取りこぼしジッタが復活。
